% !TEX root = ../../tree_notes.tex
% =====================================================================
%  6.2  Build from Degrees, Sequences and Codes       sec:c25-degree
%  Exports: def:c25-degseq  thm:c25-degreal  def:c25-prufer  prop:c25-prufdeg
%           lst:c25-degbld  ex:c25-deg-run  def:c25-quad  lst:c25-quad
%           prob:c25-quad  lst:c25-quad  (practice: LC 210)
% =====================================================================
\section{Build from Degrees, Sequences and Codes}\label{sec:c25-degree}

\begin{previouslybox}
\begin{enumerate}[nosep, leftmargin=1.6em]
  \item State the handshaking identity for a tree's degrees and the exact value of
        $\sum_v\degop(v)$ for the running tree $T$. (\thmref{c25-degsum})
  \item Why can a graph with a node of in-degree~$2$ never be a rooted tree?
        (\propref{c25-root})
  \item How many edges must a realisation of a degree sequence have, and how do you
        read that straight off the sequence? (\thmref{c25-cert})
  \item In the leaf-peel of a tree, why is there \emph{always} a node of degree~$1$ to remove?
        (\chref{c01})
\end{enumerate}
\end{previouslybox}

\begin{topicintro}
Section~6.1 rebuilt a tree from \emph{pairs}. Here the input is even coarser: a list of
\emph{numbers}. You are handed a \term{degree sequence} $d_1,\dots,d_n$ and asked (i) does
\emph{any} tree have exactly these degrees, and (ii) if so, \emph{build one} -- or even
\emph{count} them. The answer is shockingly clean: a list of positive integers is a tree degree
sequence \textbf{iff} it sums to $2(n-1)$ (\thmref{c25-degreal}), and the constructive proof is a
leaf-peel that is the exact inverse of the \term{Pr\"ufer code} -- the length-$(n-2)$ sequence
that bijects labelled trees with multisets, and the reason a degree array also tells you
\emph{how many} trees realise it (\propref{c25-prufdeg}; the full bijection lives in
\chref{c32}). The same ``rebuild a branching object from aggregate data'' instinct reappears in a
completely different guise: \term{constructing a quad tree} by recursively partitioning a grid
(LC 427), where the data is a matrix and the tree is a divide-and-conquer shape. And a
\emph{topological order} (LC 210) is a linear ``tree'' peeled from in-degrees -- the identical
degree-peel mechanic. One idea -- \emph{read structure back out of degrees} -- three faces.
\textbf{Prerequisites.} \thmref{c25-degsum}, \propref{c25-root} from \secref{c25-parent}. The
Pr\"ufer decoding is \emph{used} as a construction; its bijective proof and Cayley's count are
\chref{c32}. This section adds \emph{realisability, greedy construction and counting from a
degree array}, plus recursive-shape construction.

\textbf{After this section you can}
\begin{itemize}[nosep]
  \item decide in $O(n)$ whether a degree list is realisable by a tree, and say why;         \hfill\lv{3}
  \item construct \emph{an} explicit tree realising a given degree sequence by a leaf-peel;   \hfill\lv{3}--\lv{4}
  \item read the Pr\"ufer code as $\degop(v)-1$ occurrences and count realising trees;        \hfill\lv{4}--\lv{5}
  \item build a quad tree (or any recursive partition tree) from raw grid data by divide and conquer. \hfill\lv{3}
\end{itemize}
\end{topicintro}

\begin{recognitionbox}[Pattern card -- realise a tree from degrees / codes / aggregate data]
\reckey{the input is a list of \emph{numbers}, not pairs: a \code{degree[]} array, an
\code{ancestor-count}, a Pr\"ufer-like \emph{code}, a grid you must turn into a partition tree,
or an in-degree table you peel into an order.}
\recuse{
  \item ``does a tree exist with these degrees? / build one'' -- \thmref{c25-degreal} plus the leaf-peel;
  \item ``how many labelled trees have this degree sequence?'' -- Pr\"ufer multinomial \propref{c25-prufdeg};
  \item ``compress this $n\times n$ grid into a recursive tree'' -- quad tree, \probref{c25-quad};
  \item ``order the tasks so every dependency comes first'' -- in-degree peel (LC 210, the practice).
}
\recnotuse{
  \item you already have the edges and only need the root -- \secref{c25-parent};
  \item the sequence is an \emph{arbitrary graph}'s degrees (not necessarily a tree): that is the
        Erd\H{o}s--Gallai / Havel--Hakimi question for general graphs, out of scope here;
  \item you must answer subtree/path queries on a trusted tree -- \chref{c24}.
}
One-line test: \emph{is the input a pile of \textbf{numbers} that must become a
\textbf{shape}? Check the sum, then peel leaves -- this section.}
\end{recognitionbox}

\subsection{Intuition}

Forget the labels for a moment. Every tree on $n$ nodes has $n-1$ edges, and each edge donates
$1$ to two endpoints, so the degrees must sum to $2(n-1)$ -- there is simply no other total
(\thmref{c25-degsum}). The surprising direction is the \emph{converse}: as long as every $d_i\ge1$
and the sum is $2(n-1)$, you can \emph{always} wire up a tree with exactly those degrees. The
recipe is a peeling game. Because the average degree is $2(n-1)/n<2$, some node has degree~$1$ --
a forced \emph{leaf}. Pin that leaf onto any node that still has degree to spare ($>1$), spend one
unit from each, and recurse on the smaller problem. Each pin removes exactly one leaf and never
creates an illegal degree, so after $n-1$ pins you have a genuine tree. This is the same bookkeeping
that underlies Pr\"ufer encoding, seen from the decode side.

\begin{figure}[H]
\centering
\begin{tikzpicture}[diagram]
  % degree row for T (nodes 1..8)
  \node[annot, anchor=east] at (0,0) {\small $\degop$};
  \foreach \x/\val in {1/2,2/3,3/3,4/1,5/1,6/1,7/2,8/1} {
    \node[cellnode] at (\x,0) {\val};
  }
  % highlight the leaves (deg 1): 4,5,6,8
  \fill[successgreen!20] (3.62,-0.38) rectangle (4.38,0.38);
  \fill[successgreen!20] (4.62,-0.38) rectangle (5.38,0.38);
  \fill[successgreen!20] (5.62,-0.38) rectangle (6.38,0.38);
  \fill[successgreen!20] (7.62,-0.38) rectangle (8.38,0.38);
  % deg-1 row (Prufer multiplicity)
  \node[annot, anchor=east] at (0,-1.4) {\small $\degop{-}1$};
  \foreach \x/\val in {1/1,2/2,3/2,4/0,5/0,6/0,7/1,8/0} {
    \node[cellnode] at (\x,-1.4) {\val};
  }
  \node[annot, align=center] at (4.5,-2.8)
    {top row sums to $14=2(n-1)$ \checkmark\ realisable;\\
     green cells are the forced leaves the peel starts with;\\
     bottom row sums to $6=n-2$ and \emph{is} the Pr\"ufer multiset of $T$};
\end{tikzpicture}
\caption{Two arithmetic views of the same degree array. The \emph{top} row is $T$'s degrees; the
$\sum=2(n-1)$ test is the whole realisability certificate (\thmref{c25-degreal}). The \emph{bottom}
row subtracts $1$ from each entry -- its sum is $n-2$, and reading node $v$ exactly $\degop(v)-1$
times gives the Pr\"ufer \emph{multiset} $(2,2,1,3,3,7)$ of \secref{c25-parent}'s teaser.}\label{fig:c25-degrow}
\end{figure}

\subsection{Definitions and Notation}

\begin{definition}{Tree degree sequence; realisability}{c25-degseq}
A list $d=(d_1,\dots,d_n)$ of positive integers is \term{tree-realisable} if there exists a tree
$G$ on nodes $\{1,\dots,n\}$ with $\degop(i)=d_i$ for every $i$. The construction problem asks for
one such $G$; the counting problem asks how many labelled $G$ realise $d$.
\end{definition}

\noindent\textbf{In plain words.} ``Can these numbers be the degrees of \emph{some} tree, and if so
show me the tree (or count them).'' Nothing about positions or edges is given -- only the per-node
totals.

\begin{theorem}{Degree-sequence realisability for trees}{c25-degreal}
A list $d_1,\dots,d_n$ of positive integers is tree-realisable \textbf{if and only if}
$\displaystyle\sum_{i=1}^{n} d_i = 2(n-1)$. When $n\ge2$, the leaf-peel of \coderef{c25-degbld}
constructs a realising tree.
\end{theorem}
\begin{proofbox}
\emph{(Only if)} Every tree has $n-1$ edges and each edge contributes $2$ to the degree sum
(\thmref{c25-degsum}), so any realisation forces $\sum d_i=2(n-1)$, and clearly each $d_i\ge1$
(a node in a connected graph on $n\ge2$ vertices has a neighbour).
\emph{(If)} We run the peel. Since every $d_i\ge1$ and $\sum d_i=2(n-1)<2n$, the average is below
$2$, so some node has $d_i=1$ -- a leaf exists at every stage. Pin a current leaf~$\ell$ to any node
$p$ of current degree $>1$ and append edge $(\ell,p)$, decrementing both. The invariant ``all
remaining degrees positive and their sum equals twice the edges still to add'' is preserved: the sum
drops by $2$ (one from $\ell$, one from $p$) and we remove one leaf, so there are $n-1$ iterations
before all degrees reach~$0$. The $n-1$ produced edges connect all $n$ nodes (each new edge ties a
not-yet-used leaf into the growing component, and the final edge joins the last two leaves), giving
$n-1$ edges and, by \thmref{c25-cert}, a tree whose degrees are exactly $d$. \end{proofbox}
\begin{insightbox}
This is a \emph{sharp} contrast with general graphs: for an arbitrary graph the realizability test
is the Erd\H{o}s--Gallai inequalities or a Havel--Gakimi peel. For \emph{trees} the single equation
$\sum d_i=2(n-1)$ (plus $d_i\ge1$) is necessary \emph{and} sufficient -- trees are rigid enough that
one number decides everything. The degree array is ``missing information'' only about \emph{which}
tree, never about \emph{whether} one exists.
\end{insightbox}

\begin{definition}{Pr\"ufer code (construction view)}{c25-prufer}
The \term{Pr\"ufer code} of a labelled tree on $\{1,\dots,n\}$ is a sequence of length $n-2$ over
the alphabet $\{1,\dots,n\}$ obtained by repeatedly deleting the smallest leaf and recording its
neighbour. Crucially, $\degop(v)-1$ equals the number of times $v$ appears in the code, so
\emph{any} ordering with those multiplicities decodes to \emph{a} tree realising $d$. (The
deletion--recording bijection and its inverse are proven in full in \chref{c32}; here we only
\emph{use} the decode.)
\end{definition}

\begin{proposition}{Counting realising trees from a degree sequence}{c25-prufdeg}
If $d_1,\dots,d_n$ is tree-realisable ($n\ge2$), the number of \emph{labelled} trees whose degree
sequence is exactly $d$ is
\[
  \frac{(n-2)!}{\prod_{i=1}^{n}(d_i-1)!}.
\]
Summing this multinomial over all realisable $d$ recovers Cayley's $n^{\,n-2}$ labelled trees.
\end{proposition}
\begin{proofbox}
By \defref{c25-prufer}, a labelled tree is uniquely coded by a length-$(n-2)$ sequence in which
symbol $i$ appears exactly $d_i-1$ times; the sum of these multiplicities is
$\sum(d_i-1)=2(n-1)-n=n-2$, so the count is the number of distinct permutations of that multiset --
the multinomial $(n-2)!/\prod(d_i-1)!$. Since the Pr\"ufer map is a bijection (proven in
\chref{c32}), this multinomial \emph{is} the number of realising labelled trees. \end{proofbox}

\begin{example}{How many trees look like $T$'s degrees?}{c25-deg-count}
$d=[2,3,3,1,1,1,2,1]$, so $d_i-1=[1,2,2,0,0,0,1,0]$ and the multinomial is
$6!/(1!\,2!\,2!\,0!\,0!\,0!\,1!\,0!)=720/4=\mathbf{180}$ distinct labelled trees share $T$'s degree
sequence. Cayley gives $8^{6}=262144$ trees in total on $8$ labels; the $180$ above are just those
with this particular degree profile -- the array narrows the world by three orders of magnitude, yet
still cannot name \emph{which} tree. \end{example}

\subsection{The algorithm: peel leaves to realise a degree sequence}

\begin{lstlisting}[caption={Construct a tree realising a degree sequence, or report it impossible}, label={lst:c25-degbld}]
#include <bits/stdc++.h>
using namespace std;

// vertices are 0..n-1; deg[v] = required degree of v.
// Returns n-1 edges realising deg, or an empty list if deg is not tree-realisable.
vector<pair<int,int>> treeFromDegrees(const vector<int>& deg) {
    int n = (int)deg.size();
    long long sum = 0;
    for (int d : deg) { if (d < 1) return {}; sum += d; }     // positive required
    if (n == 1) return {};                                   // single node, no edges
    if (sum != 2LL * (n - 1)) return {};                     // the WHOLE test (Thm 6.2)

    vector<int> cur = deg;                                   // live degrees during the peel
    // Pruefer-style leaf-min-heap: always peel the smallest current leaf.
    priority_queue<int, vector<int>, greater<int>> leaf;
    for (int i = 0; i < n; ++i) if (cur[i] == 1) leaf.push(i);

    // Order vertices by (deg-1) copies -> a Pruefer sequence of length n-2.
    vector<int> prufer;
    for (int i = 0; i < n; ++i) for (int k = 1; k < deg[i]; ++k) prufer.push_back(i);

    vector<pair<int,int>> E;
    for (int x : prufer) {                                   // decode: attach a leaf to x
        int l = leaf.top(); leaf.pop();
        E.emplace_back(l, x);
        if (--cur[x] == 1) leaf.push(x);                     // x became a leaf
    }
    int a = leaf.top(); leaf.pop();                          // last two leaves join
    int b = leaf.top(); leaf.pop();
    E.emplace_back(a, b);
    return E;                                                // exactly n-1 edges, degrees = deg
}
\end{lstlisting}
\begin{complexitybox}
Time $\bigO{n\log n}$ -- the Pr\"ufer multiset holds $n-2$ entries and every heap operation is
$\log n$; building the multiset is itself $O(n)$ since $\sum(\degop(v)-1)=n-2$. Memory $\bigO{n}$.
The realisability \emph{decision} alone is $O(n)$ (one sum). The construction is exactly the
Pr\"ufer \emph{decoder} run on a conveniently chosen code, which is why it always succeeds when
\thmref{c25-degreal}'s equation holds.
\end{complexitybox}

\subsection{Worked Example}

\begin{example}{Realising $T$'s degree sequence by the peel}{c25-deg-run}
Feed $d=[2,3,3,1,1,1,2,1]$ (nodes $0..7$ relabel $1..8$) to \coderef{c25-degbld}. The sum is $14$
$=2(n-1)$ \checkmark, so we proceed. Initial leaves (degree~$1$): $\{4,5,6,8\}$ (1-indexed). The
Pr\"ufer multiset, node $v$ listed $\degop(v)-1$ times, is $\{1,\,2,2,\,3,3,\,7\}$ (length~$6$).
Decoding in ascending order: attach leaf $4\!\to\!1$ (node~$1$ drops to~$1$, joins the leaf pool);
attach leaf $5\!\to\!2$; attach leaf $6\!\to\!2$ (node~$2$ now degree~$1$, a leaf); attach leaf
$1\!\to\!3$; attach leaf $2\!\to\!3$ (node~$3$ drops to~$1$); then node~$7$; the final two free
leaves join. The result is \emph{some} tree with degrees exactly $[2,3,3,1,1,1,2,1]$ -- not
necessarily $T$ itself, but a legal realisation, which is all the theorem promises. Re-running with a
different Pr\"ufer \emph{order} would yield a different one of the $180$ trees counted in
\exref{c25-deg-count}.
\end{example}
\begin{tracebox}[The peel never gets stuck]
At every step the ``positive degrees, sum $=$ twice-remaining-edges'' invariant holds, so the
average degree of the live graph is below $2$ and a degree-$1$ node is guaranteed to be in the heap.
If you \emph{do} hit an empty leaf heap, the input failed \thmref{c25-degreal} -- e.g.\ $d=[2,2,2]$
sums to $6\ne4$, or $d=[3,1]$ has a node wanting $3$ with only one partner: both return empty.
\end{tracebox}

\subsection{From grids to recursive trees: the quad tree}

The degree array is aggregate data of type ``counts''. A \term{quad tree} is aggregate data of type
``a square matrix'', and the reconstruction is \emph{recursive partition}: if a square is uniform,
it collapses to a leaf; otherwise split it into four quadrants and rebuild each. The output tree's
shape is forced by the grid -- no edges are ever given, yet a genuine tree pops out. This is the
construction soul of LC 427.

\begin{figure}[H]
\centering
\begin{tikzpicture}[diagram, scale=0.75, transform shape]
  % ---- 4x4 grid on the left ----
  \begin{scope}[xshift=-3.2cm]
    \node[annot] at (1,-0.4) {\small input grid};
    \foreach \i in {0,1,2,3} \foreach \j in {0,1,2,3} {
      \draw[draw=primaryblue!60, line width=0.6pt] (\i,-\j) rectangle (\i+1,-\j-1);
    }
    % shade the 1-cells: bottom-left 2x2 all 1, plus top-left single
    \fill[accentorange!35] (0,-2) rectangle (2,-4);   % bottom-left quadrant = 1
    \node[font=\scriptsize] at (0.5,-2.5) {$1$}; \node[font=\scriptsize] at (1.5,-2.5) {$1$};
    \node[font=\scriptsize] at (0.5,-3.5) {$1$}; \node[font=\scriptsize] at (1.5,-3.5) {$1$};
    \node[font=\scriptsize] at (0.5,-0.5) {$0$}; \node[font=\scriptsize] at (1.5,-0.5) {$0$};
    \node[font=\scriptsize] at (2.5,-0.5) {$0$}; \node[font=\scriptsize] at (3.5,-0.5) {$0$};
    \node[font=\scriptsize] at (0.5,-1.5) {$0$}; \node[font=\scriptsize] at (1.5,-1.5) {$0$};
    \node[font=\scriptsize] at (2.5,-1.5) {$0$}; \node[font=\scriptsize] at (3.5,-1.5) {$0$};
    \node[font=\scriptsize] at (2.5,-2.5) {$1$}; \node[font=\scriptsize] at (3.5,-2.5) {$0$};
    \node[font=\scriptsize] at (2.5,-3.5) {$0$}; \node[font=\scriptsize] at (3.5,-3.5) {$0$};
    \draw[draw=warningred, line width=1pt, dashed] (2,-0) -- (2,-4);
    \draw[draw=warningred, line width=1pt, dashed] (0,-2) -- (4,-2);
  \end{scope}
  % ---- quad tree on the right ----
  \node[troot] (r) at (0,0) {\scriptsize root (split)};
  \node[tleaf] (tl) at (-2.6,-1.6) {\scriptsize leaf $0$};
  \node[tleaf] (tr) at (-0.9,-1.6) {\scriptsize leaf $0$};
  \node[tnode] (bl) at (0.9,-1.6) {\scriptsize leaf $1$};
  \node[tnode] (br) at (2.6,-1.6) {\scriptsize split};
  \node[tleaf] (b1) at (1.6,-3) {\scriptsize $0$};
  \node[tleaf] (b2) at (2.4,-3) {\scriptsize $0$};
  \node[tleaf] (b3) at (3.2,-3) {\scriptsize $0$};
  \node[tleaf] (b4) at (4.0,-3) {\scriptsize $0$};
  \draw[tedge] (r) -- (tl); \draw[tedge] (r) -- (tr);
  \draw[tedge] (r) -- (bl); \draw[tedge] (r) -- (br);
  \draw[tedge] (br) -- (b1); \draw[tedge] (br) -- (b2);
  \draw[tedge] (br) -- (b3); \draw[tedge] (br) -- (b4);
\end{tikzpicture}
\caption{The same grid, two representations. The $4\times4$ grid on the left is split (dashed) into
four $2\times2$ quadrants; three are uniform (leaf $0$, leaf $0$, leaf $1$) and the bottom-right
still mixes $0$/$1$, so it recurses one more level. The quad tree on the right is the
\emph{reconstruction}: a node is a leaf exactly when its square is uniform, internal otherwise --
LC 427.}\label{fig:c25-quad}
\end{figure}

\begin{definition}{Region quad tree}{c25-quaddef}
A \term{quad tree} over an $n\times n$ grid is a tree whose root covers the whole square. A node
covering a square $S$ is a \term{leaf} if every cell of $S$ holds the same value (that value is
stored on the leaf); otherwise it is \term{internal} and has exactly four children covering the
top-left, top-right, bottom-left and bottom-right halves of $S$.
\end{definition}

\subsection{Variations and Twists}

\begin{itemize}[nosep]
  \item \textbf{Sorted vs labelled degrees}: \thmref{c25-degreal} is about the labelled model (node
        $i$ has degree $d_i$). If only the \emph{multiset} of degrees matters, divide the count by
        automorphisms -- still a Pr\"ufer argument (\chref{c32}). \hfill\lv{4}--\lv{5}
  \item \textbf{Bounded-degree trees}: constrain which $d_i$ are allowed (e.g.\ binary: all $d_i\le3$
        for a rooted binary tree); realisability then needs a matching argument, not just the sum. \hfill\lv{5}
  \item \textbf{Prefer a specific tree}: to \emph{rebuild $T$ itself}, the degree array is not enough --
        you need a full Pr\"ufer sequence (the order matters), \chref{c32}. \hfill\lv{4}
  \item \textbf{Quad tree, faster uniformity test}: replace the $O(\text{area})$ scan with a 2-D prefix
        sum of the values so ``is $S$ uniform?'' is $O(1)$; total build $\bigO{n^{2}}$. \hfill\lv{4}
  \item \textbf{k-d tree / octree}: same partition-tree idea with $k$ dimensions or splits on
        a data point instead of uniformity -- construction-by-recursion, LC 427's template. \hfill\lv{3}--\lv{4}
\end{itemize}

\subsection{Common Pitfalls}\label{sec:c25-degree-pitfalls}

\begin{warningbox}[Trap 1 -- accepting a sequence with a zero or a bad sum]
$d_i=0$ with $n\ge2$ is impossible in a connected graph (an isolated node), and $\sum d_i\ne2(n-1)$
is fatal -- yet both sneak past code that only ``starts peeling''. Test positivity \emph{and} the sum
\emph{before} constructing (see the two early \code{return \{\}} lines of \coderef{c25-degbld}).
\end{warningbox}
\begin{lstlisting}[caption={\textcolor{warningred}{WRONG}: peel without validating}, language=C++, basicstyle=\ttfamily\small]
for (int i=0;i<n;i++) if (deg[i]==1) leaf.push(i);   // sum may be != 2(n-1): heap can empty mid-run
\end{lstlisting}
\begin{lstlisting}[caption={\textcolor{successgreen}{RIGHT}: check the certificate first}, language=C++, basicstyle=\ttfamily\small]
long long s=0; for(int d:deg){ if(d<1) return {}; s+=d; }
if (s != 2LL*(n-1)) return {};                        // reject before any peel
\end{lstlisting}

\begin{warningbox}[Trap 2 -- assuming the constructed tree is \emph{the} tree]
A degree sequence fixes only the \emph{degrees}, not the wiring. Unless a Pr\"ufer sequence's
\emph{order} is given, you recover \emph{one of} the $\prod$-many realisations -- fine for
``build any'', wrong for ``recover exactly $T$''. Do not claim the answer is unique when
\propref{c25-prufdeg} returns a count $>1$.
\end{warningbox}

\begin{warningbox}[Trap 3 -- quad tree: splitting uniform squares]
The base case is \emph{uniform value}, not \emph{size~$1$}. A $2\times2$ block of all $1$s is a single
leaf; recursing past a uniform square spawns empty children and corrupts \code{isLeaf}. Test
uniformity \emph{first}, and only split when it fails (\coderef{c25-quad}).
\end{warningbox}

\subsection{Problem Lab}

The anchor, LC 427, is the purest ``data $\to$ tree'' reconstruction in the chapter: the edges are
never mentioned -- they are \emph{implied} by the recursive partition of a grid.

% ---------------------------------------------------------------------------
\solvedproblem{c25-quad}{Construct Quad Tree}{LeetCode 427}
  {https://leetcode.com/problems/construct-quad-tree/}{LC Medium L3}
  {divide and conquer, recursive construction, uniformity test, tree building from grid}

\begin{problemstatement}[Construct Quad Tree]
Given an $n\times n$ grid \code{grid} of $0$s and $1$s, build its \textbf{quad tree} and return the
root. A \textbf{leaf} has \code{isLeaf=true} and \code{val} equal to the (single) value in its
region. An \textbf{internal} node has \code{isLeaf=false}, arbitrary \code{val}, and four children
\code{topLeft}, \code{topRight}, \code{bottomLeft}, \code{bottomRight} covering the four quadrants.
$n$ is a power of~$2$, $1\le n\le 6$.
\end{problemstatement}

\begin{attemptgate}[Attempt it yourself first -- \lv{3}, give yourself 25 minutes]
Fair first attempt: recursion on a square. Base case when the whole square is one value.
\peek{How do you decide ``the square $(r_0,c_0)\!\to\!(r_1,c_1)$ is uniform''? And when it
\emph{isn't}, into how many sub-squares does it break, and what are their corners?}
\nudge{Scan the square; if all cells equal $\mathit{val}$, return \code{new Node(val, true)}.
Otherwise split at the mid row/col into four rectangles -- top-left $[r_0,rm)\times[c_0,cm)$, etc.\
-- and recurse; return \code{new Node(val, false, tl, tr, bl, br)}.}
\fullhint{\code{build(r0,c0,r1,c1)}: let $v=\mathit{grid}[r_0][c_0]$; if the region is all $v$
return a leaf; else $rm=(r_0{+}r_1)/2$, $cm=(c_0{+}c_1)/2$ and recurse on the four half-open
quadrants, threading them as the four children. Top call \code{build(0,0,n,n)}. Full build and
proof below.}
\end{attemptgate}

\solstep{Understanding the Problem}
The tree's shape is completely determined by the grid: a node is a leaf iff its region is uniform.
The four children are ordered, so this is a \emph{fixed} 4-ary recursion -- the only real work is
(a) testing uniformity and (b) getting the quadrant boundaries right (half-open rectangles so the
four pieces tile the parent exactly). Unlike \secref{c25-parent} there are no edges to read -- the
edges are \emph{created} by the partition.

\solstep{Key Observations}
\begin{observation}
$n$ is a power of two, so halving always lands on integer boundaries and bottoms out at $1\times1$
(a trivially uniform square), guaranteeing the recursion terminates.
\end{observation}
\begin{observation}
``Uniform'' is a property of a region, checked independently at every node; the naive scan is
$O(\text{area})$. Because LC 427 caps $n\le6$, the naive total $O(n^2\log n)$ is trivially fine --
but a 2-D prefix sum makes each test $O(1)$ and the whole build $O(n^2)$ for large grids.
\end{observation}
\begin{observation}
For an internal node the \code{val} is unspecified by the problem, so any consistent choice works;
the standard is the top-left cell's value -- it must not be mistaken for a leaf decision.
\end{observation}

\solstep{Think It Through With Me}
\begin{thinkbox}
\deadend{Build the tree by first converting the grid into an explicit edge list.}
{There \emph{are} no edges in the input -- you would invent a graph, then re-derive the very
partition you were trying to encode. Quad trees are built \emph{top-down by recursion}, not from pairs.}
\deadend{Split before testing uniformity.}{Then a solid $4\times4$ block of $1$s becomes $16$ tiny
leaves and four levels of pointless children -- you produce a non-minimal quad tree and fail the
\code{isLeaf} checks. Uniformity is the \emph{base case}; test it first.}
\aha{A node is a leaf \emph{iff} its square is uniform. So: test, then either return a leaf or split
into the four half-open quadrants and recurse -- the recursion \emph{is} the construction.}
\end{thinkbox}

\solstep{Solution}
Define $\mathit{build}(r_0,c_0,r_1,c_1)$ over the half-open square $[r_0,r_1)\times[c_0,c_1)$. Read
$v=\mathit{grid}[r_0][c_0]$; if every cell in the square equals $v$, return a leaf \code{Node(v,true)}.
Otherwise compute midpoints $rm,cm$ and return an internal \code{Node(v,false,\ldots)} built from
the four recursive calls on the top-left $[r_0,rm)\times[c_0,cm)$, top-right $[r_0,rm)\times[cm,c_1)$,
bottom-left $[rm,r_1)\times[c_0,cm)$ and bottom-right $[rm,r_1)\times[cm,c_1)$ quadrants. The root
is $\mathit{build}(0,0,n,n)$.

\solstep{Proof of Correctness}
\begin{proofbox}
\emph{Termination}: each recursive call halves the side length, reaching $1\times1$ (uniform,
hence a leaf), so the depth is $\log_2 n$ and every branch terminates. \emph{Leaf correctness}:
$\mathit{build}$ returns \code{isLeaf=true} only after scanning the whole region and confirming all
cells equal $v$, matching the definition. \emph{Internal correctness}: when the uniformity scan
fails, the four half-open quadrants are disjoint and their union is exactly the parent square, so the
four children cover the parent's region with no gaps or overlaps, each built recursively to
represent its sub-region. By induction on region size, every returned node faithfully represents its
square, and the root represents the full grid -- so the constructed quad tree is the required one. \end{proofbox}

\solstep{Implementation}
\begin{lstlisting}[caption={Construct Quad Tree -- recursive partition of a uniform test (LeetCode 427)}, label={lst:c25-quad}]
#include <bits/stdc++.h>
using namespace std;

// Definition for a quad-tree node.
class Node {
public:
    bool val;
    bool isLeaf;
    Node* topLeft;
    Node* topRight;
    Node* bottomLeft;
    Node* bottomRight;
    Node() : val(false), isLeaf(false),
             topLeft(nullptr), topRight(nullptr),
             bottomLeft(nullptr), bottomRight(nullptr) {}
    Node(bool _val, bool _isLeaf) : val(_val), isLeaf(_isLeaf),
             topLeft(nullptr), topRight(nullptr),
             bottomLeft(nullptr), bottomRight(nullptr) {}
    Node(bool _val, bool _isLeaf, Node* a, Node* b, Node* c, Node* d)
        : val(_val), isLeaf(_isLeaf), topLeft(a), topRight(b), bottomLeft(c), bottomRight(d) {}
};

class Solution {
    vector<vector<int>> g;
    bool uniform(int r0, int c0, int r1, int c1, int v) {
        for (int i = r0; i < r1; ++i)
            for (int j = c0; j < c1; ++j)
                if (g[i][j] != v) return false;
        return true;
    }
    Node* build(int r0, int c0, int r1, int c1) {
        int v = g[r0][c0];
        if (uniform(r0, c0, r1, c1, v)) return new Node(v == 1, true);   // leaf (base case)
        int rm = (r0 + r1) / 2, cm = (c0 + c1) / 2;                     // split into 4
        return new Node(v == 1, false,
            build(r0, c0, rm, cm), build(r0, cm, rm, c1),              // top-left, top-right
            build(rm, c0, r1, cm), build(rm, cm, r1, c1));             // bottom-left, bottom-right
    }
public:
    Node* construct(vector<vector<int>>& grid) {
        g = grid;
        int n = (int)grid.size();
        return build(0, 0, n, n);
    }
};
\end{lstlisting}

\solstep{Trace}
\begin{tracebox}[The \figref{c25-quad} grid]
\figref{c25-quad}: bottom-left $2\times2$ is all~$1$; top-left and top-right $2\times2$ are all~$0$;
bottom-right is mixed. \code{build(0,0,4,4)} scans the whole $4\times4$ -- mixed $\Rightarrow$ internal.
Four calls: \code{build(0,0,2,2)} uniform $0\Rightarrow$ leaf; \code{build(0,2,2,4)} uniform $0\Rightarrow$
leaf; \code{build(2,0,4,2)} uniform $1\Rightarrow$ leaf; \code{build(2,2,4,4)} mixed $\Rightarrow$ split
into four $1\times1$ leaves $(1,0,0,0)$. The returned tree has root with $3$ leaf children and one
internal child with $4$ leaves -- exactly \figref{c25-quad}. \checkmark
\end{tracebox}

\solstep{Complexity}
\begin{complexitybox}
Time $O(n^2\log n)$ with the naive scan: level $0$ scans $n^2$, the $4$ children together scan
$4\cdot(n/2)^2=n^2$, and so on for $\log n$ levels -- $O(n^2)$ per level. Space $O(n^2)$ nodes in the
worst checkerboard (every $1\times1$ is its own leaf, $O(n^2)$ nodes). With $n\le6$ this is free; for
large grids, a 2-D prefix sum makes \code{uniform} $O(1)$ and the build $O(n^2)$. Verdict: \ac.
\end{complexitybox}

\solstep{Pitfalls and Variants}
\begin{warningbox}
Three breakers: (i) splitting a uniform region (Trap~3) -- test \emph{first}. (ii) Off-by-one
quadrants: use half-open $[r_0,r_1)$ and midpoints so the four children tile exactly; closed
boundaries overlap and one cell gets two parents. (iii) Forgetting \code{isLeaf} on internal nodes
(\code{new Node(v,false,\dots)}) -- the judge compares structure, so a mis-flagged leaf fails.
\emph{Variant:} LC 427's sibling is the ``intersection of two quad trees'' and ``quad-tree
compression'' -- same recursion, extra combine step.
\end{warningbox}

\begin{insightbox}[Key takeaways]
\begin{enumerate}[nosep, leftmargin=1.6em]
  \item \textbf{Degrees are a certificate.} One equation, $\sum d_i=2(n-1)$, decides existence
        outright (\thmref{c25-degreal}) -- far cleaner than the general-graph case.
  \item \textbf{The peel \emph{is} Pr\"ufer decoding.} Leaf-peel construction and Pr\"ufer counting
        are the same mechanism seen constructively vs. enumeratively (\chref{c32}).
  \item \textbf{Recursive partition is reconstruction too.} A quad tree turns grid data into a tree
        with no edges ever named -- construction from aggregate data, one shape per data value.
\end{enumerate}
\end{insightbox}

\begin{hardeningbox}
\begin{itemize}[nosep, leftmargin=1.4em]
  \item \textbf{O(1) uniformity}: build a 2-D prefix sum \code{pre} of the grid; a region is uniform
        iff its cell-sum equals $0$ or equals its area -- turns \probref{c25-quad} into $O(n^2)$.
  \item \textbf{Recover \emph{exactly} $T$}: replace \coderef{c25-degbld}'s chosen Pr\"ufer order with
        the given code, then decode -- the standard Prüfer$\to$tree decoder (\chref{c32}).
  \item \textbf{Existential edge count}: among the $180$ realisations of $T$'s degrees, ask how many
        have node $2$ adjacent to node $7$; a Pr\"ufer-window counting exercise, \chref{c32}.
\end{itemize}
\end{hardeningbox}

\subsection{Practice Problems}

\begin{stuckbox}
\begin{itemize}[nosep, leftmargin=1.4em]
  \item Did you validate positivity \emph{and} the $\sum=2(n-1)$ equation before peeling?
        (Trap~1, \thmref{c25-degreal}.)
  \item Did you resist claiming uniqueness when \propref{c25-prufdeg} gives a count $>1$? (Trap~2.)
  \item In the quad tree, did you test uniformity \emph{before} splitting? (Trap~3, \probref{c25-quad}.)
  \item Are your quadrant rectangles half-open so the four children tile the parent? (\probref{c25-quad}.)
\end{itemize}
\end{stuckbox}

\begin{practiceset}
  \practice{Course Schedule II}{LeetCode 210}
    {https://leetcode.com/problems/course-schedule-ii/}{LC Medium L3}
    {Given prerequisite \emph{pairs}, return a valid ordering of all courses, or \code{[]}. This is the
    \textbf{in-degree peel} of \thmref{c25-degreal} on a \emph{directed acyclic} graph: repeatedly emit
    every node of current in-degree~$0$ (a ``source'', the DAG analogue of a degree-$1$ leaf), delete
    its outgoing edges, and repeat (Kahn's algorithm). Recognise the shared mechanic -- peel by degree,
    reconstruct an order/structure from aggregate counts -- and note the \emph{cycle detector}: if the
    peel finishes with nodes still having positive in-degree, the input is not a DAG (the ``realisability
    fails'' moment). Rebuild \code{adj} from the pairs first, exactly \secref{c25-parent}.}
\end{practiceset}

\begin{recallbox}[10-second refresher -- 6.2 Degrees, Sequences \& Codes]
A positive list is a tree degree sequence \emph{iff} $\sum d_i=2(n-1)$ (\thmref{c25-degreal}); build a
realisation by peeling the smallest leaf onto a degree-$>1$ node (\coderef{c25-degbld}) -- the same
move as Pr\"ufer decoding, where node $v$ appears $\degop(v)-1$ times (\defref{c25-prufer}). The count
of realising labelled trees is $(n-2)!/\prod(d_i-1)!$ (\propref{c25-prufdeg}); full bijection and
Cayley are \chref{c32}. A quad tree reconstructs a grid by uniform-then-split recursion (\probref{c25-quad});
a topological order reconstructs a DAG by the in-degree peel (LC 210).
\selftest{
  \item Decide realisability of $[3,2,2,1,1,1]$ for $n=6$: yes or no, and by which test?
  \item Run one peel step on $T$'s degrees -- which node do you pin first and to what?
  \item How many labelled trees share a degree array with $\degop-1$ multiset $\{1,2,2,1\}$?
  \item In LC 427, why must uniformity be tested \emph{before} you split a square?
  \item Explain how Kahn's algorithm is the directed in-degree sibling of the leaf-peel.
}
\end{recallbox}
